
\newgeometry{hmargin=1cm,vmargin=1.2cm}
\begin{landscape}
\setlength{\parskip}{0pt}%
\vspace*{\fill}
\centering
\begin{tikzpicture}[x=1cm, y=0.86cm,
  place/.style={circle, draw=black!80, line width=0.7pt,
    minimum size=8.5mm, inner sep=0pt},
  trans/.style={rectangle, draw=black!80, line width=0.7pt, fill=white,
    minimum width=9.5mm, minimum height=4.5mm, inner sep=1.5pt,
    font=\footnotesize},
  every label/.style={font=\footnotesize, label distance=1.5pt},
  arc/.style={line width=0.55pt, draw=black!80,
    arrows={Bar[width=2.1mm]}-{Stealth[length=2.4mm]}},
  carc/.style={arc, draw=czero},
  rarc/.style={line width=0.55pt, draw=black!80,
    arrows={Stealth[length=2.4mm]}-{Stealth[length=2.4mm]}},
  elide/.style={line width=0.55pt, draw=czero, dotted,
    -{Stealth[length=2.4mm]}},
  cell/.style={draw=hndeep, fill=hnaccent!50, rounded corners=1pt,
    font=\footnotesize, inner sep=3pt, minimum height=4.2mm},
  fold/.style={draw=hndeep, fill=hnaccent!50, rounded corners=1pt,
    font=\footnotesize, inner sep=3pt, minimum height=4.2mm,
    align=center},
  defcell/.style={cell, dashed},
  plate/.style={draw=hnmid, line width=0.7pt, rounded corners=2pt,
    fill=hnaccent!18, inner sep=5pt},
  platelab/.style={font=\tiny\ttfamily, text=hnmid, inner sep=1pt},
  wire/.style={-, line width=0.5pt, draw=black!75},
  wlab/.style={font=\scriptsize, inner sep=1pt},
  port/.style={font=\small\itshape, inner sep=1.5pt},
  ann/.style={font=\scriptsize, text=black!70, align=left,
    anchor=north west, inner sep=0pt}]

\node[ann] at (-0.2,1.55) {\emph{the plant of the Petri page, elided
  to the diffusion re-entrancy (steps 1 and 5 share the machine
  pool); everything ungated is omitted}};

\node[place, label=above:{$q_1$}] (q1) at (1.0,0) {};
\node[trans] (s1) at (3.0,0) {$\mathrm{start}_1$};
\node[place, draw=wkA, line width=0.9pt, label=above:{$p_1$}] (p1) at (4.8,0) {};
\node[trans] (f1) at (6.6,0) {$\mathrm{finish}_1$};
\node[place, label=above:{$q_5$}] (q5) at (12.6,0) {};
\node[trans] (s5) at (14.4,0) {$\mathrm{start}_5$};
\node[place, draw=wkE, line width=0.9pt, label=above:{$p_5$}] (p5) at (16.2,0) {};
\node[trans] (f5) at (18.0,0) {$\mathrm{finish}_5$};
\node[place, label=above:{$q_6$}] (q6) at (19.8,0) {};

\draw[arc] (q1) -- (s1);
\draw[arc, draw=wkA] (s1) -- (p1);
\draw[carc] (p1) -- (f1);
\draw[elide] (f1) to[bend left=14]
  node[wlab, above, text=czero] {steps 2--4 (implantation, lithography) --- ungated} (q5);
\draw[arc] (q5) -- (s5);
\draw[arc, draw=wkE] (s5) -- (p5);
\draw[carc] (p5) -- (f5);
\draw[arc] (f5) -- (q6);
\draw[elide] (q6.east) -- ++(1.45,0)
  node[wlab, right, text=czero, align=left]
  {step 6, ship\\(the arrival gate of the Petri page)};

\node[place, label={[font=\footnotesize]180:{$m_{\mathrm{diff}}$}}]
  (mdiff) at (8.7,-1.05) {};
\fill[black] ([xshift=-1.2mm]mdiff.center) circle (0.065);
\fill[black] ([xshift= 1.2mm]mdiff.center) circle (0.065);
\draw[arc, rounded corners=5pt] (mdiff.170) -| ([xshift=-2mm]s1.south);
\draw[arc, rounded corners=5pt] (mdiff.10)  -| ([xshift=-2mm]s5.south);
\draw[arc, rounded corners=5pt] ([xshift=2mm]f1.south) |- (mdiff.140);
\draw[arc, rounded corners=5pt] ([xshift=2mm]f5.south) -- ++(0,-0.62)
  -| (mdiff.40);

\begin{pgfonlayer}{holebg}
  \draw[dashed, line width=0.7pt, draw=black!75]
        (-0.2,-12.1) rectangle (24.6,-2.05);
\end{pgfonlayer}
\node[font=\footnotesize\itshape] at (12.2,-11.8)
  {Supervisor seat --- filled by the synthesized policy $\pi^{\ast}$};

\node[place, fill=white,
  label={[font=\footnotesize]160:{$g_1$}}] (g1) at (3.0,-2.05) {};
\fill[black] (g1.center) circle (0.065);
\node[place, fill=white,
  label={[font=\footnotesize]160:{$g_5$}}] (g5) at (14.4,-2.05) {};
\draw[rarc, draw=wkA] (g1) -- ([xshift=2mm]s1.south);
\draw[rarc, draw=wkE] (g5) -- ([xshift=2mm]s5.south);
\node[place, fill=white,
  label={[font=\footnotesize]200:{$c_1$}}] (c1) at (6.6,-2.05) {};
\node[place, fill=white,
  label={[font=\footnotesize]340:{$c_5$}}] (c5) at (18.0,-2.05) {};
\draw[arc, draw=wkA] (f1.south) -- (c1.north);
\draw[arc, draw=wkE] ([xshift=-2mm]f5.south) -- ++(0,-0.35) -| (c5.north);
\node[rectangle, draw=black!80, line width=0.7pt, fill=white,
  minimum size=3.5mm, inner sep=0pt,
  label={[font=\footnotesize]200:{$q_5$ pool}}] (tq5) at (12.0,-2.05) {};
\node[rectangle, draw=black!80, line width=0.7pt, fill=white,
  minimum size=3.5mm, inner sep=0pt,
  label={[font=\footnotesize]200:{$p_5$ pool}}] (tp5) at (16.8,-2.05) {};
\draw[rarc] (q5.south west) -- (tq5.north);
\draw[rarc] (p5.south east) -- (tp5.north);

\begin{pgfonlayer}{boxbg}
  \draw[line width=0.9pt, draw=black!80, fill=white]
        (1.2,-11.35) rectangle (24.2,-3.35);
\end{pgfonlayer}
\node[font=\footnotesize\bfseries, anchor=north east, align=left,
  text width=5.35cm] at (24.05,-3.55)
  {$\pi^{\ast}$ --- one transition, fired every tick: the kernel
   reads the four boundary observables and rewrites the two gates};

\draw[rarc] (c1.south) -- (c1.south |- 0,-3.35);
\draw[rarc] (c5.south) -- (c5.south |- 0,-3.35);
\draw[rarc] (tq5.south) -- (tq5.south |- 0,-3.35);
\draw[rarc] (tp5.south) -- (tp5.south |- 0,-3.35);
\draw[arc, draw=wkA] ([xshift=-1.2mm]g1.south) -- ++(0,-1.0);
\draw[arc, draw=wkA] ([xshift=1.2mm]g1.south) ++(0,-1.0) -- ++(0,1.0);
\draw[arc, draw=wkE] ([xshift=-1.2mm]g5.south) -- ++(0,-1.0);
\draw[arc, draw=wkE] ([xshift=1.2mm]g5.south) ++(0,-1.0) -- ++(0,1.0);

\node[port] (Pg1) at (3.0,-3.75) {$g_1$};
\node[port] (Pc1) at (6.6,-3.75) {$c_1$};
\node[port] (Pq5) at (12.0,-3.75) {$q_5$};
\node[port] (Pg5) at (14.4,-3.75) {$g_5$};
\node[port] (Pp5) at (16.8,-3.75) {$p_5$};
\node[port] (Pc5) at (18.0,-3.75) {$c_5$};

\node[cell] (sc1) at (11.2,-5.7) {$1$};
\node[fold] (sf1) at (13.0,-5.7) {fold $\Sigma$};
\node[cell] (scw) at (16.9,-5.7) {$w$};
\node[fold] (sfw) at (16.9,-7.3) {fold $\Sigma$};
\begin{pgfonlayer}{plates}
  \node[plate, fit=(sc1)] (pl1) {};
  \node[plate, fit=(scw)] (plw) {};
\end{pgfonlayer}
\node[platelab, anchor=north east] at (pl1.south east) {plate $q_5$};
\node[platelab, anchor=north east] at (plw.south east) {plate $p_5$};
\draw[wire, rounded corners=4pt] (Pq5.south) -- ++(0,-0.55)
  -- node[wlab, above left, pos=0.6] {$[q_5]$} (sc1.north);
\draw[wire, rounded corners=4pt] (Pp5.south)
  -- node[wlab, right, pos=0.45] {$[p_5]$} (scw.north);
\draw[wire] (sc1.east) -- (sf1.west);
\draw[wire] (scw.south) -- (sfw.north);

\node[defcell] (thc) at (21.6,-7.3) {margin $\theta(c_5)$};
\node[defcell] (phc) at (7.0,-8.6) {quota $\varphi(c_1, c_5)$};
\draw[wire, rounded corners=4pt] (Pc5.south) -- (18.0,-7.3)
  -- (thc.west);
\draw[wire, rounded corners=4pt] (Pc1.south) -- (6.6,-7.0)
  -- node[wlab, left, pos=0.85] {$c_1$} (6.6,-8.35);
\draw[wire, rounded corners=4pt] (18.0,-7.3) -- (18.0,-9.6)
  -- node[wlab, above, pos=0.65] {$c_5$} (7.0,-9.6) -- (phc.south);

\node[cell] (G5) at (15.4,-8.6)
  {$\ell \ge 1 \;\wedge\; u \le \theta$};
\node[cell] (G1) at (3.9,-6.6) {$c_1 < \varphi$};
\draw[wire, rounded corners=4pt] (sf1.east) -- ++(0.35,0)
  -- node[wlab, left, pos=0.55] {$\ell$} (14.05,-8.35)
  -- (G5.west |- 0,-8.5) -- (G5.west);
\draw[wire, rounded corners=4pt] (sfw.south)
  -- node[wlab, right, pos=0.35] {$u$} (16.9,-8.0)
  -- (16.1,-8.0) -- (16.1,-8.25);
\draw[wire, rounded corners=4pt] (thc.south)
  |- node[wlab, above, pos=0.78] {$\theta$} (G5.east);
\draw[wire, rounded corners=4pt] (phc.west)
  -- node[wlab, above, pos=0.6] {$\varphi$} (3.9,-8.6) -- (G1.south);
\draw[wire, rounded corners=4pt] (6.6,-5.2) -- (3.9,-5.2) -- (G1.north);
\draw[wire, rounded corners=4pt] (G5.north) -- (15.4,-4.6)
  -- node[wlab, above, pos=0.35] {open/hold} (14.4,-4.6) -- (Pg5.south);
\draw[wire, rounded corners=4pt] (G1.west) -- (3.0,-6.6)
  -- node[wlab, left, pos=0.6] {open/hold} (Pg1.south);

\node[ann, text width=21.6cm] at (1.7,-10.15)
  {closed forms of the two dashed defined-symbol cells ---
   $\theta$ (diffusion-slot margin): the first $1 + \min(N-2,\,2)$
   step-5 slots run consecutively on one machine --- $\theta = 0$, the
   two-tick handover; every later slot alternates machines,
   pre-positioned one tick early --- $\theta = 1$.\quad
   $\varphi$ (feed quota): unbounded while $c_1 < \min(N,\,4)$
   (step-1 services back-to-back); afterwards one admission per slot
   window (services 5, 6 into the first freed windows, then
   alternating).};
\end{tikzpicture}%
\par\medskip
\begin{minipage}{0.94\linewidth}
\platecaption{The synthesized policy $\pi^{\ast}$ (optimal for $N \le 7$), drawn
as what it is: a filler for the supervisor seat. The dashed seat of
Figure~\ref{fig:minifab-sdcpn-petri} is here nearly the whole page,
and the plant survives only as the gated neighbourhood along the top
--- the diffusion re-entrancy, where $\mathrm{start}_1$ and
$\mathrm{start}_5$ compete for the same two machines. Straddling the
boundary, the seat's declared interface: the two pooled reads, the two
firing counters, and the two gate places --- a strict subset of the
fifteen-observable, four-act full-dispatch interface of
section~\ref{sec:minifab-fulldispatch}, and the only wires
$\pi^{\ast}$ needs. Inside sits one enormous Petri transition: the
policy's firing kernel, applied once per tick, reading the four
observables and rewriting the two gates. Its interior is a hypernet in the plate notation of
Figure~\ref{fig:minifab-cert-hypernets}: the pooled observables are
two plated folds ($\ell = \Sigma\,1$ over $[q_5]$, the queue length;
$u = \Sigma\,w$ over $[p_5]$, the in-service residual ---
pre-positioning keeps at most the incumbent in $p_5$, so the sum reads
its residual), the counters decode the schedule phase through the two
\emph{dashed} cells ($\theta$, $\varphi$ --- defined symbols, not
primitive generators: piecewise-constant functions of the counters,
closed forms in the box), and two threshold cells (region-membership)
emit the gate bits. Holding $\mathrm{start}_5$ until $u \le \theta$
pre-positions each diffusion slot; holding $\mathrm{start}_1$ at the
quota keeps a diffusion machine free for step~5 --- deliberate idling,
the one capability the never-idle class lacks. What this filler
attains, each claim at its tier (section~\ref{sec:minifab-optimum}):
the $N$-th arrival into the last queue at $255N + 371 + d(N)$
ticks ($d=0,0,2$, then $4$) for $N=1,\dots,12$ and $N=84$.
The weekly replay reaches $21{,}795$ against the decoded
$N=84$ container plus a transcription of the rows it does not
seal; the $N=1,\dots,12$ replays use that model with the
pending-lot count and unlock threshold substituted. Optimality
on the simplified net is established for $N=1,\dots,7$
and is open beyond; the
Lean machine-contiguity floor is for an abstract schedule model. The best tested never-idle policy needs
$255N + 455$ at $N = 84$ --- these two gates are worth $80$ ticks over it
(whether an untested work-conserving policy closes that gap is open), and replaying $\pi^{\ast}$ with
holds disallowed is refused at its first delayed diffusion start.
Under deterministic closed dynamics, imprecise-decision-theoretic
selection provably collapses to makespan minimization (Lean-checked)
--- so where $\pi^{\ast}$ is optimal ($N \le 7$), this page is also
the decision-theoretic optimum. Initial
marking drawn: $g_1$ open ($c_1 = 0 < \varphi$), $g_5$ shut
($\ell = 0$).}{fig:minifab-pistar}
\end{minipage}
\par
\vspace*{\fill}
\thispagestyle{empty}
\end{landscape}
\restoregeometry
